The belief that a silent passage of flatus smells worse than an audible one appears in many languages but has not
been tested. We follow a passage from sound through dispersion to a bystander's memory, using a lip-reed model of the
anal valve, analytic and computational dispersion models and literature observer models. Most releases, 85\%, never
set the valve oscillating. We prove that under linear transport a slower release of the same gas cannot give a higher
peak, in any room or flow. A silent release can win only near the source, where a loud jet is pre-diluted by a
factor set by its exit-to-ambient velocity ratio, the Silencer number. In computational fluid dynamics of seated and
standing people, the warm air rising along the body takes a slow release and lets a fast one pass, so momentum
decides which nose a release reaches. As emitted, a silent passage is the stronger one at a bystander's nose in
\gap{v4} of pairs. A silent passage is noticed only if smelled, and this selection alone can explain the belief: a
remembered silent passage seems the stronger one in \gap{v4} of comparisons. The explanation predicts that releasers
hold the belief less firmly than bystanders and that anosmic people never remember a silent passage.
